%% 2-12 文字式のやくそく（12ページ、p.150〜161）　組5
%% 第11回：やくそくごとに「なぜ」を書く。数を入れて確かめる段・図・例題を足した。消しゴムは共通の \bkeraser。
%% ×をはぶく・÷は分数・x×x=x^2、項と係数、同類項＝同じ種類（単位）どうししかたせない、
%% 分配法則でかっこをはずす、まとめてわるときは1つずつわってよい（5-2・5-4 の思い出そう）、指数のきまりを文字で（5-4 の思い出そう）
\makeatletter
% 長さ x のテープ（文字入り）  \bkiimtape{(x,y)}{長さmm}{色}{文字}
\newcommand\bkiimtape[4]{\begin{scope}[shift={#1}]
  \filldraw[fill=#3!22, draw=#3!80!black, line width=0.8pt, rounded corners=0.6mm] (0,0) rectangle (#2,6.5);
  \node[font=\small, inner sep=0pt] at (#2/2,3.25) {#4};\end{scope}}
% 文字式を、まん中に1行で置く  \bkiimline{…}
\newcommand\bkiimline[1]{\par\addvspace{3pt}{\centering #1\par}\addvspace{3pt}}
\makeatother

\bkchapter{2}{12}{\R{文字式}{も|じ|しき}のやくそく}{\bklevel{2}}{\R{文字式}{も|じ|しき}のやくそく}
\bkmichishirube{2-2,2-7,2-11}{2-13,2-14,3-8,5-2,5-4}

\begin{goal}
\begin{itemize}[label={\color{bkTheme}●}, leftmargin=1.6\zw, itemsep=0pt, topsep=1pt]
\item \bklevel{2}\enspace ×をはぶき、÷を\R{分数}{ぶん|すう}にして、\R{文字式}{も|じ|しき}を\R{短}{みじか}く\R{書}{か}ける。
\item \bklevel{2}\enspace \R{同類項}{どう|るい|こう}をまとめたり、かっこをはずしたりできる。
\item \bklevel{2}\enspace \R{指数}{し|すう}のきまりを、\R{文字}{も|じ}を\R{使}{つか}って\R{書}{か}いて\R{使}{つか}える。
\end{itemize}
\end{goal}

\begin{hatena}
えんぴつ3\R{本}{ぼん}とえんぴつ2\R{本}{ほん}を\R{合}{あ}わせると、えんぴつ5\R{本}{ほん}。では、えんぴつ3\R{本}{ぼん}とけしゴム2こを\R{合}{あ}わせると、「えんぴつけしゴム5こ」かな？
\begin{bkfigure}
\begin{tikzpicture}[x=1mm, y=1mm]
  \foreach \i in {0,1,2} {\bkpencil{(0,\i*5)}{24}{bkBlue}}
  \node[font=\large] at (30,5) {＋};
  \foreach \i in {0,1} {\bkpencil{(35,\i*5+2.5)}{24}{bkBlue}}
  \node[font=\large] at (65,5) {＝};
  \node[anchor=west, font=\sffamily\bfseries\small] at (70,5) {えんぴつ5\R{本}{ほん}};
  \begin{scope}[shift={(0,-22)}]
  \foreach \i in {0,1,2} {\bkpencil{(0,\i*5)}{24}{bkBlue}}
  \node[font=\large] at (30,5) {＋};
  \bkeraser[bkBlue]{(35,2.5)}{11.5}{5}\bkeraser[bkBlue]{(49,2.5)}{11.5}{5}
  \node[font=\large] at (65,5) {＝};
  \node[anchor=west, font=\sffamily\bfseries\small, text=bkRed] at (70,5) {えんぴつけしゴム5こ？};
  \end{scope}
\end{tikzpicture}
\end{bkfigure}
\R{長}{なが}さでも\R{同}{おな}じです。\mbox{\qty{3}{m}＋\qty{2}{m}＝\qty{5}{m}}。では、\mbox{\qty{3}{m}＋2\R{秒}{びょう}}は？　\R{文字式}{も|じ|しき}の $3\times x+2\times y$ は、$5\times x\times y$ にまとめてよいのでしょうか？
\end{hatena}

この\R{章}{しょう}では、\R{文字式}{も|じ|しき}を\R{短}{みじか}く\R{書}{か}くやくそくと、\R{文字式}{も|じ|しき}をまとめたり、かっこをはずしたりするきまりを\R{学}{まな}びます。

\clearpage
%% ---------------- やってみよう・思い出そう ----------------
\begin{yattemiyou}[\R{電卓}{でん|たく}で、2つの\R{計算}{けい|さん}をくらべよう]
\begin{enumerate}[label={\bkcircled[bkGreen]{\arabic*}}, leftmargin=2.4\zw, itemsep=1pt, topsep=2pt]
\item $x$ に2、3、1.5を\R{入}{い}れて、$x^2\times x^3$ と $x^5$ と $x^6$ を\R{電卓}{でん|たく}で\R{計算}{けい|さん}しよう。
\item $A$、$B$、$h$ に\R{表}{ひょう}の\R{数}{かず}を\R{入}{い}れて、$(A+B)\div h$ と $A\div h+B\div h$ を\R{計算}{けい|さん}しよう。
\item いつも\R{同}{おな}じになるのは、どれとどれかな？
\end{enumerate}
\begin{bkcenter}
\begin{bktable}{colspec={Q[c,m,wd=8mm] Q[c,m,wd=22mm] Q[c,m,wd=10mm] Q[c,m,wd=10mm]}, rowsep=1pt, colsep=4pt, cells={font=\small}, row{1}={bg=bkGreen!14}}
$x$ & $x^2\times x^3$ & $x^5$ & $x^6$ \\
2 & $4\times8=32$ & 32 & 64 \\
3 & & & \\
1.5 & & & \\
\end{bktable}\par\vspace{7pt}
\begin{bktable}{colspec={*{3}{Q[c,m,wd=6mm]} Q[c,m,wd=24mm] Q[c,m,wd=28mm]}, rowsep=1pt, colsep=4pt, cells={font=\small}, row{1}={bg=bkGreen!14}}
$A$ & $B$ & $h$ & $(A+B)\div h$ & $A\div h+B\div h$ \\
6 & 4 & 2 & $10\div2=5$ & $3+2=5$ \\
9 & 3 & 3 & & \\
5 & 7 & 4 & & \\
\end{bktable}
\end{bkcenter}
\end{yattemiyou}

\begin{omoidasou}
\bkrecall{2-7}{$10^2\times10^3=10^5$（かける\R{回数}{かい|すう}をたす）。$10^0=1$、$10^{-1}=\dfrac{1}{10}$}
\bkrecall{2-11}{\R{分配法則}{ぶん|ぱい|ほう|そく} $(a+b)\times c=a\times c+b\times c$}
\bkrecall{2-3}{たし\R{算}{ざん}・ひき\R{算}{ざん}ができるのは、\R{同}{おな}じ\R{種類}{しゅ|るい}の\R{量}{りょう}どうしだけ。}
\bkrecall{2-4}{わり\R{算}{ざん}は\R{分数}{ぶん|すう}で\R{書}{か}ける：$2\div3=\dfrac{2}{3}$}
\bkrecall{2-2}{$(-1)\times3=-3$、$(-1)\times(-2)=2$。×$(-1)$ で\R{向}{む}きが\R{反対}{はん|たい}になる。}
\end{omoidasou}


\clearpage
%% ---------------- 数学の道具 ----------------
\bksection{doguu}{\R{数学}{すう|がく}の\R{道具}{どう|ぐ}}

\bksubsection{1}{×をはぶいて、\R{短}{みじか}く\R{書}{か}く}

\leavevmode\label{bk:2-12:yakusoku}2-11では、$x=v\times t$ のように×を\R{書}{か}きました。\R{文字式}{も|じ|しき}には×がたくさん\R{出}{で}てくるので、\R{数学}{すう|がく}や\R{物理}{ぶつ|り}では、\R{次}{つぎ}のやくそくで\R{短}{みじか}く\R{書}{か}きます。

\begin{bkcenter}
\begin{bktable}{colspec={Q[c,m,wd=6mm] Q[l,m,wd=56mm] Q[c,m,wd=30mm] Q[c,m,wd=22mm]}, rowsep=1pt, cells={font=\small}, row{1}={font=\sffamily\bfseries\small, bg=bkTheme!12}}
 & やくそく & もとの\R{式}{しき} & \R{短}{みじか}く\R{書}{か}くと \\
1 & \R{文字}{も|じ}のかけ\R{算}{ざん}の×をはぶく & $v\times t$、\ $3\times x$ & $vt$、\ $3x$ \\
2 & \R{数}{かず}は\R{文字}{も|じ}の\R{前}{まえ}に\R{書}{か}く & $x\times3$ & $3x$ \\
3 & 1は\R{書}{か}かない & $1\times x$、\ $(-1)\times x$ & $x$、\ $-x$ \\
4 & \R{同}{おな}じ\R{文字}{も|じ}のかけ\R{算}{ざん}は\R{累乗}{るい|じょう}で（2-7） & $x\times x$、\ $x\times x\times x$ & $x^2$、\ $x^3$ \\
5 & ÷は\R{分数}{ぶん|すう}で\R{書}{か}く & $x\div3$、\ $a\div b$ & $\dfrac{x}{3}$、\ $\dfrac{a}{b}$ \\
6 & \R{文字}{も|じ}はアルファベット\R{順}{じゅん}に & $y\times x\times2$ & $2xy$ \\
\end{bktable}
\bktabcaption{\R{文字式}{も|じ|しき}のやくそく}
\end{bkcenter}

\noindent\begin{minipage}[c]{0.5\linewidth}
$3x$ は $3\times x$ で、「3が $x$ こ\R{分}{ぶん}」とも「$x$ が3つ\R{分}{ぶん}」とも\R{読}{よ}めます（\R{交換法則}{こう|かん|ほう|そく}、2-11）。\R{長}{なが}さ $x$ のテープ3\R{本}{ぼん}\R{分}{ぶん}の\R{長}{なが}さです。\R{数}{かず}と\R{数}{かず}のあいだの×は、はぶきません。$3\times4$ を $34$ と\R{書}{か}くと、\R{三十四}{さん|じゅう|よん}とまちがえるからです。
\end{minipage}\hfill
\begin{minipage}[c]{0.46\linewidth}\centering
\begin{tikzpicture}[x=1mm, y=1mm]
  \foreach \i in {0,1,2} {\bkiimtape{(\i*18,0)}{18}{bkBlue}{$x$}}
  \bkbrace{(0,-1)}{(54,-1)}{$3x$}{-1}
\end{tikzpicture}
\bkfigcaption{$3x=x+x+x$}
\end{minipage}

\vspace{3pt}
\textbf{それぞれのやくそくの、わけ}
\begin{enumerate}[label={\bkcircled[bkTheme]{\arabic*}}, leftmargin=2.4\zw, itemsep=1pt, topsep=2pt]
\item ×を\R{書}{か}かなくても、\R{文字}{も|じ}は1\R{字}{じ}ずつなので、$vt$ が「$v$ と $t$ のかけ\R{算}{ざん}」だとわかります。
\item \R{数}{かず}を\R{前}{まえ}にそろえると、$3x$ と $x3$ のように\R{同}{おな}じものが2\R{通}{とお}りに\R{書}{か}かれず、くらべやすくなります。
\item 1をかけても\R{大}{おお}きさは\R{変}{か}わらないので（$1\times5=5$）、$1x$ は $x$ と\R{同}{おな}じです。$(-1)\times x$ は「$x$ の\R{反対}{はん|たい}の\R{向}{む}き」なので（2-2）、$-x$ と\R{書}{か}きます。
\item $x\times x$ は、1\R{辺}{ぺん}が $x$ の\R{正方形}{せい|ほう|けい}の\R{面積}{めん|せき}です（\R{下}{した}の\R{図}{ず}）。2-7で $3\times3=3^2$ と\R{書}{か}いたのと\R{同}{おな}じです。
\item わり\R{算}{ざん}は\R{分数}{ぶん|すう}で\R{書}{か}けました（2-4：$2\div3=\dfrac{2}{3}$）。÷よりも、どこからどこまでをわるのかが、はっきりします。
\item \R{文字}{も|じ}の\R{順番}{じゅん|ばん}も、だれが\R{書}{か}いても\R{同}{おな}じ\R{形}{かたち}になるようにそろえます（$2xy$ と $2yx$ は\R{同}{おな}じ）。
\end{enumerate}

\begin{bkfigure}
\begin{tikzpicture}[x=1mm, y=1mm]
  \fill[bkBlue!18] (0,0) rectangle (20,20);
  \draw[bkBlue!70!black, line width=0.8pt] (0,0) rectangle (20,20);
  \node[font=\small, inner sep=0pt] at (10,10) {$x^2$};
  \node[anchor=north, font=\small, inner sep=0pt] at (10,-1.8) {$x$};
  \node[anchor=east, font=\small, inner sep=0pt] at (-1.8,10) {$x$};
  \node[anchor=west, font=\small, inner sep=0pt] at (26,13) {$x=3$ なら $3\times3=9$};
  \node[anchor=west, font=\small, inner sep=0pt] at (26,6) {$x=5$ なら $5\times5=25$};
\end{tikzpicture}
\bkfigcaption{$x\times x=x^2$ は、1\R{辺}{ぺん} $x$ の\R{正方形}{せい|ほう|けい}の\R{面積}{めん|せき}。}
\end{bkfigure}

\R{短}{みじか}く\R{書}{か}いても、\R{意味}{い|み}は\R{変}{か}わりません。たとえば $x=3$ のとき、$4x$ は $4\times3=12$、$x^2$ は $3\times3=9$、$\dfrac{x}{3}$ は $3\div3=1$ です。2-11の $x=v\times t$ は、これからは $x=vt$ と\R{書}{か}きます。

\bksubsection{2}{\R{同}{おな}じ\R{種類}{しゅ|るい}どうしだけ、まとめられる}

\leavevmode\label{bk:2-12:dourui}\bkhatenaref を\R{考}{かんが}えます。えんぴつ1\R{本}{ぽん}の\R{長}{なが}さを $x$ cm、けしゴム1この\R{長}{なが}さを $y$ cm とします。

\begin{bkfigure}
\begin{tikzpicture}[x=1mm, y=1mm]
  \node[anchor=east, font=\sffamily\small] at (-2,11) {$3x+2x$};
  \foreach \i in {0,1,2} {\bkiimtape{(\i*16,8)}{16}{bkBlue}{$x$}}
  \foreach \i in {3,4} {\bkiimtape{(\i*16,8)}{16}{bkBlue}{$x$}}
  \node[anchor=west, font=\sffamily\bfseries\small] at (84,11) {＝ $5x$};
  \node[anchor=east, font=\sffamily\small] at (-2,1) {$3x+2y$};
  \foreach \i in {0,1,2} {\bkiimtape{(\i*16,-2)}{16}{bkBlue}{$x$}}
  \foreach \i in {0,1} {\bkiimtape{(48+\i*9,-2)}{9}{bkOrange}{$y$}}
  \node[anchor=west, font=\sffamily\bfseries\small] at (70,1) {まとめられない};
\end{tikzpicture}
\bkfigcaption{$x$ どうしはまとめられる。$x$ と $y$ は、まとめられない。}
\end{bkfigure}

$3x+2x$ は、$x$ が3つと2つで5つなので、$5x$ です。\R{分配法則}{ぶん|ぱい|ほう|そく}（2-11）を\R{使}{つか}うと、$3x+2x=(3+2)x=5x$ です。でも、$3x+2y$ は、$x$ と $y$ の\R{長}{なが}さがちがうので、これ\R{以上}{い|じょう}まとめられません。\bkhatenaref の\R{答}{こた}えは「まとめられない」です。\qty{3}{m}と2\R{秒}{びょう}をたせないのと、\R{同}{おな}じです（2-3）。

$3x+2y$ のように、＋でつながった1つ1つ（$3x$ と $2y$）を\textbf{\R{項}{こう}}といい、\bkidx{こう}{\R{項}{こう}}\R{文字}{も|じ}にかけてある\R{数}{かず}を\textbf{\R{係数}{けい|すう}}といいます（$3x$ の\R{係数}{けい|すう}は3）。\bkidx{けいすう}{\R{係数}{けい|すう}}\R{文字}{も|じ}の\R{部分}{ぶ|ぶん}が\R{同}{おな}じ\R{項}{こう}を\textbf{\R{同類項}{どう|るい|こう}}といいます。\bkidx{どうるいこう}{\R{同類項}{どう|るい|こう}}まとめられるのは、\R{同類項}{どう|るい|こう}どうしだけです。ひき\R{算}{ざん}は、マイナスの\R{数}{かず}をたすと\R{考}{かんが}えます（2-2）。$5a-3a=5a+(-3a)$ なので、\R{項}{こう}は $5a$ と $-3a$ です。

\noindent\begin{minipage}[c]{0.52\linewidth}
ひき\R{算}{ざん}も、テープで\R{考}{かんが}えるとわかります。$5a-3a$ は、$a$ のテープ5\R{本}{ほん}から3\R{本}{ぼん}をとるので、2\R{本}{ほん}が\R{残}{のこ}ります。\R{分配法則}{ぶん|ぱい|ほう|そく}でも、$5a-3a=(5-3)a=2a$ です。
\end{minipage}\hfill
\begin{minipage}[c]{0.45\linewidth}\centering
\begin{tikzpicture}[x=1mm, y=1mm]
  \foreach \i in {0,...,4} {\bkiimtape{(\i*10,0)}{10}{bkGreen}{$a$}}
  \foreach \i in {2,3,4} {\draw[bkRed, line width=1pt] (\i*10+1.5,-1) -- (\i*10+8.5,7.5);}
  \bkbrace{(0,-1.5)}{(20,-1.5)}{$2a$}{-1}
  \node[anchor=south, font=\sffamily\footnotesize, text=bkRed, inner sep=0pt] at (35,8.6) {3\R{本}{ぼん}とる};
\end{tikzpicture}
\bkfigcaption{$5a-3a=2a$}
\end{minipage}

\vspace{3pt}
とる\R{数}{かず}のほうが\R{多}{おお}いときは、\R{係数}{けい|すう}がマイナスになります。
\[ 3x-7x=(3-7)x=-4x \]
$3-7=-4$ は、2-2の\R{数直線}{すう|ちょく|せん}で3から7\R{左}{ひだり}へ\R{進}{すす}んだ\R{所}{ところ}です。

\R{同類項}{どう|るい|こう}が\R{離}{はな}れているときは、\R{項}{こう}を\R{符号}{ふ|ごう}ごと\R{並}{なら}べかえてから、まとめます。たとえば、$4x+3+x+5=4x+x+3+5=5x+8$（$x$ の\R{項}{こう}どうし、\R{数}{かず}どうし）。$x=2$ を\R{入}{い}れて\R{確}{たし}かめると、\R{左}{ひだり}は $8+3+2+5=18$、\R{右}{みぎ}は $10+8=18$ で\R{同}{おな}じです。

$x+x^2$ は、\R{長}{なが}さ（m）と\R{面積}{めん|せき}（\unit{m^2}）のように\R{種類}{しゅ|るい}がちがうので、まとめられません。$x=3$ なら $x+x^2=3+9=12$ ですが、$2x^2=18$ や $2x=6$ にはなりません。

\bksubsection{3}{かっこをはずす・まとめてわる}

\leavevmode\label{bk:2-12:kakko}\textbf{かっこをはずす}\quad \R{分配法則}{ぶん|ぱい|ほう|そく}（2-11）を、\R{文字式}{も|じ|しき}で\R{使}{つか}います。

\noindent\begin{minipage}[c]{0.55\linewidth}
\bkiimline{$3(x+2)=3\times x+3\times2=3x+6$}
\R{右}{みぎ}の\R{図}{ず}の\R{長方形}{ちょう|ほう|けい}の\R{面積}{めん|せき}を、2\R{通}{とお}りに\R{表}{あらわ}した\R{式}{しき}です。
\end{minipage}\hfill
\begin{minipage}[c]{0.42\linewidth}\centering
\begin{tikzpicture}[x=1mm, y=1mm]
  \fill[bkBlue!20] (0,0) rectangle (30,18); \fill[bkOrange!25] (30,0) rectangle (44,18);
  \draw[bkInk, line width=0.9pt] (0,0) rectangle (44,18); \draw[bkInk, line width=0.9pt] (30,0) -- (30,18);
  \node[font=\small] at (15,9) {$3x$}; \node[font=\small] at (37,9) {$6$};
  \bkbrace{(0,19.5)}{(30,19.5)}{$x$}{1}\bkbrace{(30,19.5)}{(44,19.5)}{$2$}{1}
  \node[anchor=east, font=\small] at (-1.5,9) {$3$};
\end{tikzpicture}
\bkfigcaption{$3(x+2)=3x+6$}
\end{minipage}

\vspace{3pt}
かっこの\R{前}{まえ}がマイナスのときは、$-1$ をかけると\R{考}{かんが}えます（2-2）。
\[ -(x-2)=(-1)\times x+(-1)\times(-2)=-x+2 \]
$x=5$ で\R{確}{たし}かめます。\R{左}{ひだり}は $-(5-2)=-3$、\R{右}{みぎ}は $-5+2=-3$ で、\R{同}{おな}じです。もし $-x-2$ とまちがえると $-5-2=-7$ になり、\R{合}{あ}いません。


\textbf{まとめてわる}\quad \bkyattemiyouref の②では、$(A+B)\div h$ と $A\div h+B\div h$ が、いつも\R{同}{おな}じになりました。\R{長方形}{ちょう|ほう|けい}で\R{考}{かんが}えると、わけがわかります。

\noindent\begin{minipage}[c]{0.44\linewidth}\centering
\begin{tikzpicture}[x=1mm, y=1mm]
  \fill[bkBlue!20] (0,0) rectangle (27,15); \fill[bkOrange!25] (27,0) rectangle (45,15);
  \draw[bkInk, line width=0.9pt] (0,0) rectangle (45,15); \draw[bkInk, line width=0.9pt] (27,0) -- (27,15);
  \node[font=\small] at (13.5,7.5) {\R{面積}{めん|せき} $A$}; \node[font=\small] at (36,7.5) {\R{面積}{めん|せき} $B$};
  \node[anchor=east, font=\small] at (-1.5,7.5) {$h$};
  \bkbrace{(0,-1)}{(27,-1)}{$A\div h$}{-1}\bkbrace{(27,-1)}{(45,-1)}{$B\div h$}{-1}
  \bkbrace{(0,16.5)}{(45,16.5)}{$(A+B)\div h$}{1}
\end{tikzpicture}
\end{minipage}\hfill
\begin{minipage}[c]{0.53\linewidth}
たての\R{長}{なが}さが $h$ の\R{長方形}{ちょう|ほう|けい}の\R{横}{よこ}の\R{長}{なが}さは、\mbox{\R{面積}{めん|せき}÷$h$}。\R{全体}{ぜん|たい}の\R{横}{よこ}は $(A+B)\div h$、2つに\R{分}{わ}けると $A\div h$ と $B\div h$ です。どちらも\R{同}{おな}じ\R{横}{よこ}の\R{長}{なが}さです。
\end{minipage}
\begin{bkcenter}
{\setlength\fboxsep{3pt}\fcolorbox{bkTheme}{white}{まとめてわるときは、1つずつわってよい：$(A+B)\div h=A\div h+B\div h$}}
\end{bkcenter}
\R{分数}{ぶん|すう}で\R{書}{か}くと $\dfrac{A+B}{h}=\dfrac{A}{h}+\dfrac{B}{h}$ です。たとえば、$\dfrac{6x+4}{2}=\dfrac{6x}{2}+\dfrac{4}{2}=3x+2$。ただし、\R{分母}{ぶん|ぼ}が\R{和}{わ}の $\dfrac{6}{2+1}$ は、\R{分}{わ}けられません。\R{正}{ただ}しくは $\dfrac{6}{3}=2$ ですが、\R{分}{わ}けて $\dfrac{6}{2}+\dfrac{6}{1}=3+6=9$ とすると\R{合}{あ}いません。1つずつわってよいのは、わられる\R{数}{かず}（\R{分子}{ぶん|し}）がたし\R{算}{ざん}のときだけです。

\bksubsection{4}{\R{指数}{し|すう}のきまりを、\R{文字}{も|じ}で\R{書}{か}く}

\leavevmode\label{bk:2-12:shisu}2-7では、$10^2\times10^3=10^5$ のように、かける\R{回数}{かい|すう}をたせばよいことを、\R{数}{かず}で\R{学}{まな}びました。\bkyattemiyouref の①でも、$x$ がどんな\R{数}{かず}でも、$x^2\times x^3$ と $x^5$ は\R{同}{おな}じでした。わけは、かけている $x$ の\R{数}{かず}を\R{数}{かず}えるとわかります。

\begin{bkfigure}
\begin{tikzpicture}[x=1mm, y=1mm]
  \node[font=\small] at (-12,3.25) {$x^2\times x^3=$};
  \foreach \i in {0,1} {\bkiimtape{(\i*13,0)}{11}{bkBlue}{$x$}}
  \node[font=\small] at (26.5,3.25) {$\times$};
  \foreach \i in {0,1,2} {\bkiimtape{(31+\i*13,0)}{11}{bkOrange}{$x$}}
  \node[anchor=west, font=\small] at (72,3.25) {$=x^5$（$x$ が $2+3=5$ こ）};
\end{tikzpicture}
\bkfigcaption{$x$ を2\R{回}{かい}かけて、さらに3\R{回}{かい}かけると、5\R{回}{かい}かけたことになる。}
\end{bkfigure}

\bkyattemiyouref の\R{表}{ひょう}のように、$x=2$ なら $2^2\times2^3=4\times8=32$、$2^5=32$ で\R{同}{おな}じです。$x^6$（64）にはなりません。だから、かける\R{回数}{かい|すう}を\R{表}{あらわ}す1、2、3、…のどんな\R{数}{かず} $m$、$n$ でも、
\[ x^m\times x^n=x^{m+n} \]
ここで $m$、$n$ は「\R{指数}{し|すう}に\R{入}{はい}る\R{数}{かず}の\R{箱}{はこ}」です（2-11）。$x$ そのものは、$x$ を1\R{回}{かい}だけ\R{使}{つか}うので $x=x^1$ と\R{考}{かんが}えます。
わり\R{算}{ざん}では、\R{分母}{ぶん|ぼ}と\R{分子}{ぶん|し}で\R{同}{おな}じ $x$ を\R{消}{け}して、$\dfrac{x^5}{x^2}=\dfrac{x\times x\times x\times x\times x}{x\times x}=x^3$。かける\R{回数}{かい|すう}をひきます。

\R{次}{つぎ}に、2-7の\R{表}{ひょう}のならびを\R{文字}{も|じ}で\R{書}{か}きます。\R{右}{みぎ}へ1つ\R{進}{すす}むたびに、$x$ でわります。0ではわれないので（1-14）、ここでは $x$ は0でない\R{数}{かず}とします。

\begin{bkcenter}
\begin{bktable}{colspec={Q[l,m,wd=16mm] *{6}{Q[c,m,wd=15mm]}}, rowsep=2pt, cells={font=\small}}
 & $x^3$ & $x^2$ & $x^1$ & $x^0$ & $x^{-1}$ & $x^{-2}$ \\
\SetCell{font=\small} \R{意味}{い|み} & $x\times x\times x$ & $x\times x$ & $x$ & 1 & $\dfrac{1}{x}$ & $\dfrac{1}{x^2}$ \\
$x=2$ & 8 & 4 & 2 & 1 & $\dfrac{1}{2}$ & $\dfrac{1}{4}$ \\
\end{bktable}
\bktabcaption{\R{右}{みぎ}へ1つ\R{進}{すす}むたびに $x$ でわる}
\end{bkcenter}

$x^1\div x=x\div x=1$（\R{同}{おな}じ\R{数}{かず}どうしのわり\R{算}{ざん}）なので $x^0=1$。$x$ が2でも10でも、0\R{乗}{じょう}は1です。もう1\R{回}{かい}わって $x^{-1}=\dfrac{1}{x}$、もう1\R{回}{かい}わって $x^{-2}=\dfrac{1}{x^2}$ です。マイナスの\R{指数}{し|すう}でも、$x^m\times x^n=x^{m+n}$ が\R{使}{つか}えます。たとえば、$x^3\times x^{-1}=x^2$。これは、$x^3\div x=x^2$ と\R{同}{おな}じです。

\begin{motto}[2]
\textbf{$x^0=1$ を、もう1つの\R{方法}{ほう|ほう}で\R{確}{たし}かめる}\par
わり\R{算}{ざん}のきまり「かける\R{回数}{かい|すう}をひく」を、\R{分母}{ぶん|ぼ}と\R{分子}{ぶん|し}が\R{同}{おな}じ $\dfrac{x^2}{x^2}$ に\R{使}{つか}ってみます。
\[ \frac{x^2}{x^2}=x^{2-2}=x^0 \]
いっぽう、\R{同}{おな}じ\R{数}{かず}どうしのわり\R{算}{ざん}なので、$\dfrac{x^2}{x^2}=\dfrac{x\times x}{x\times x}=1$ です。\R{同}{おな}じものを2\R{通}{とお}りに\R{計算}{けい|さん}したので、$x^0=1$ と\R{決}{き}めれば、きまりがいつでも\R{使}{つか}えます。\par
\R{同}{おな}じように、$\dfrac{x^2}{x^3}=x^{2-3}=x^{-1}$ で、\R{約分}{やく|ぶん}すると $\dfrac{x\times x}{x\times x\times x}=\dfrac{1}{x}$。だから $x^{-1}=\dfrac{1}{x}$ です。\R{表}{ひょう}のならびで\R{考}{かんが}えても、\R{約分}{やく|ぶん}で\R{考}{かんが}えても、\R{同}{おな}じ\R{答}{こた}えになります。
\end{motto}


\begin{taisetsu}
×をはぶく（$3\times x=3x$）、÷は\R{分数}{ぶん|すう}（$x\div3=\dfrac{x}{3}$）、$x\times x=x^2$\par\vspace{2pt}
\R{同類項}{どう|るい|こう}（\R{同}{おな}じ\R{種類}{しゅ|るい}）どうしだけ、まとめられる：$3x+2x=5x$\par\vspace{2pt}
まとめてわるときは、1つずつわってよい：$(A+B)\div h=A\div h+B\div h$\par\vspace{2pt}
$x^m\times x^n=x^{m+n}$、\ $\dfrac{x^m}{x^n}=x^{m-n}$\par\vspace{2pt}
$x^0=1$、\ $x^{-1}=\dfrac{1}{x}$、\ $x^{-2}=\dfrac{1}{x^2}$
\end{taisetsu}

\begin{kotoba}
\small
\bkword{\R{項}{こう}}{＋でつながった1つ1つ。$3x+2y$ の\R{項}{こう}は $3x$ と $2y$。}
\bkword{\R{係数}{けい|すう}}{\R{文字}{も|じ}にかけてある\R{数}{かず}。$3x$ の\R{係数}{けい|すう}は3、$-x$ の\R{係数}{けい|すう}は $-1$。}
\bkword{\R{同類項}{どう|るい|こう}}{\R{文字}{も|じ}の\R{部分}{ぶ|ぶん}が\R{同}{おな}じ\R{項}{こう}。$3x$ と $2x$ は\R{同類項}{どう|るい|こう}。$x$ と $x^2$ はちがう。}
\end{kotoba}

\bkkigocards{\bkkigocard{$x^{-1}$}{エックスのマイナス1\R{乗}{じょう}}{$\dfrac{1}{x}$ のこと。$x$ で1\R{回}{かい}わる。}{x to the minus one}{2-12}}
{\bkkigocard{$x^0$}{エックスの0\R{乗}{じょう}}{1のこと（$x$ は0でない）。}{x to the zero}{2-12}}

%% ---------------- 例題 ----------------
\clearpage
\bksection{reidai}{\R{例題}{れい|だい}でたしかめよう}

\begin{reidai}[やくそくで\R{短}{みじか}く\R{書}{か}く]
\R{短}{みじか}く\R{書}{か}きましょう。\quad（1）\nobreak $a\times5\times b$\quad（2）\nobreak $x\times x\times(-3)$\quad（3）\nobreak $(x+1)\div4$\par
\kangaekata \pageref{bk:2-12:yakusoku}ページの\R{表}{ひょう}のやくそくを、1つずつあてはめます。\par
\shiki（1）\R{数}{かず}を\R{前}{まえ}に、×をはぶく。$a\times5\times b=5ab$\par
\hspace*{2.6\zw}（2）$x\times x=x^2$、\R{数}{かず}を\R{前}{まえ}に。$x\times x\times(-3)=-3x^2$\par
\hspace*{2.6\zw}（3）÷は\R{分数}{ぶん|すう}。かっこの\R{中}{なか}がまとまって\R{分子}{ぶん|し}になる。$\dfrac{x+1}{4}$\par
\kotae（1）\nobreak $5ab$\quad（2）\nobreak $-3x^2$\quad（3）\nobreak $\dfrac{x+1}{4}$
\end{reidai}

\begin{jibun}
\R{短}{みじか}く\R{書}{か}きましょう。\quad（1）\nobreak $y\times x\times2$\quad（2）\nobreak $a\times a\times a$\quad（3）\nobreak $(a-b)\div3$
\bkjans{（1）\nobreak $2xy$\quad（2）\nobreak $a^3$\quad（3）\nobreak $\dfrac{a-b}{3}$}
\end{jibun}

\begin{reidai}[\R{同類項}{どう|るい|こう}をまとめる]
まとめましょう。\quad（1）\nobreak $4x+3x$\quad（2）\nobreak $5a+2-3a+4$\quad（3）\nobreak $2x+3x^2$\quad（4）\nobreak $3x+1-7x$\par
\kangaekata \R{同類項}{どう|るい|こう}どうし、\R{数}{かず}の\R{項}{こう}どうしをまとめます。\par
\shiki（1）$4x+3x=(4+3)x=7x$\par
\hspace*{2.6\zw}（2）$5a+2-3a+4=5a-3a+2+4=2a+6$\par
\hspace*{2.6\zw}（3）$x$ と $x^2$ は\R{同類項}{どう|るい|こう}ではないので、まとめられない\par
\hspace*{2.6\zw}（4）\R{符号}{ふ|ごう}ごと\R{並}{なら}べかえる。$3x-7x+1=(3-7)x+1=-4x+1$\par
\kotae（1）\nobreak $7x$\quad（2）\nobreak $2a+6$\quad（3）\nobreak $2x+3x^2$ のまま\quad（4）\nobreak $-4x+1$
\end{reidai}

\begin{jibun}
まとめましょう。\quad（1）\nobreak $7y-2y$\quad（2）\nobreak $3x-5+x+8$\quad（3）\nobreak $a^2+a$\quad（4）\nobreak $2b+6-5b$
\bkjans{（1）\nobreak $5y$\quad（2）\nobreak $4x+3$\quad（3）\nobreak まとめられない（$a^2+a$ のまま）\quad（4）\nobreak $-3b+6$}
\end{jibun}

\begin{reidai}[かっこをはずす・まとめてわる]
\R{計算}{けい|さん}しましょう。\quad（1）\nobreak $3(2x+1)$\quad（2）\nobreak $-(4a-3)$\quad（3）\nobreak $(8x+6)\div2$\par
\kangaekata かっこの\R{外}{そと}の\R{数}{かず}を、\R{中}{なか}のどの\R{項}{こう}にもかけます。\par
\shiki（1）$3(2x+1)=3\times2x+3\times1=6x+3$\par
\hspace*{2.6\zw}（2）$-(4a-3)=(-1)\times4a+(-1)\times(-3)=-4a+3$\par
\hspace*{2.6\zw}（3）$(8x+6)\div2=\dfrac{8x}{2}+\dfrac{6}{2}=4x+3$\par
\kotae（1）\nobreak $6x+3$\quad（2）\nobreak $-4a+3$\quad（3）\nobreak $4x+3$
\end{reidai}

\begin{jibun}
\R{計算}{けい|さん}しましょう。\quad（1）\nobreak $5(x-2)$\quad（2）\nobreak $-(2y+5)$\quad（3）\nobreak $(9a-6)\div3$
\bkjans{（1）\nobreak $5x-10$\quad（2）\nobreak $-2y-5$\quad（3）\nobreak $3a-2$}
\end{jibun}

\begin{motto}[2]
\textbf{\R{答}{こた}えが\R{正}{ただ}しいか、\R{数}{かず}を\R{入}{い}れて\R{確}{たし}かめよう}\par
\R{文字式}{も|じ|しき}の\R{計算}{けい|さん}は、\R{形}{かたち}を\R{変}{か}えても\R{値}{あたい}が\R{変}{か}わらないように\R{書}{か}きかえることです。だから、\R{文字}{も|じ}にかんたんな\R{数}{かず}を\R{入}{い}れて、もとの\R{式}{しき}と\R{答}{こた}えの\R{式}{しき}が\R{同}{おな}じ\R{値}{あたい}になるかを\R{見}{み}れば、まちがいに\R{気}{き}づけます。
\begin{bkcenter}
\begin{bktable}{colspec={Q[l,m,wd=38mm] Q[c,m,wd=34mm] Q[c,m,wd=34mm]}, rowsep=1.5pt, cells={font=\small}, row{1}={font=\sffamily\bfseries\small, bg=bkGray!12}}
$x=2$ を\R{入}{い}れる & もとの\R{式}{しき} $3(2x+1)$ & \R{答}{こた}え $6x+3$ \\
\R{正}{ただ}しい\R{答}{こた}え & $3\times5=15$ & $12+3=15$ \\
まちがい $6x+1$ のとき & $3\times5=15$ & $12+1=13$ \\
\end{bktable}
\end{bkcenter}
15と13で\R{合}{あ}わないので、$6x+1$ はまちがいだとわかります。1はかけても\R{大}{おお}きさが\R{変}{か}わらない\R{数}{かず}、0はかけると\R{答}{こた}えが0になる\R{数}{かず}なので、0や1を\R{入}{い}れると、まちがいを\R{見}{み}のがすことがあります。2や3を\R{使}{つか}うとよいでしょう。\R{次}{つぎ}の\R{例題}{れい|だい}でも、この\R{方法}{ほう|ほう}で\R{確}{たし}かめます。
\end{motto}

\begin{reidai}[かっこをはずしてから、まとめる]
$3(x+2)-(x-4)$ を\R{計算}{けい|さん}しましょう。$x=1$ を\R{入}{い}れて\R{確}{たし}かめましょう。\par
\kangaekata \bkcircled[bkTheme]{1}\,それぞれのかっこをはずす（\R{前}{まえ}がマイナスのかっこは、どの\R{項}{こう}にも $-1$ をかける）。\bkcircled[bkTheme]{2}\,\R{同類項}{どう|るい|こう}をまとめる。\par
\shiki \bkcircled[bkTheme]{1}\,$3(x+2)-(x-4)=3x+6-x+4$\par
\hspace*{2.6\zw}\bkcircled[bkTheme]{2}\,$3x-x+6+4=(3-1)x+10=2x+10$\par
\hspace*{2.6\zw}\R{確}{たし}かめ：$x=1$ のとき、もとの\R{式}{しき}は $3\times3-(-3)=9+3=12$\par
\hspace*{6.6\zw}$2x+10$ は $2+10=12$ で、\R{同}{おな}じ\par
\kotae $2x+10$
\end{reidai}

\begin{jibun}
$2(a+1)-(a-5)$ を\R{計算}{けい|さん}しましょう。$a=2$ を\R{入}{い}れて\R{確}{たし}かめましょう。
\bkjans{$2a+2-a+5=a+7$。$a=2$ のとき、もとの\R{式}{しき}は $6-(-3)=9$、$a+7=9$}
\end{jibun}

\begin{reidai}[\R{指数}{し|すう}のきまりを\R{使}{つか}う]
\R{計算}{けい|さん}しましょう。\quad（1）\nobreak $x^3\times x^4$\quad（2）\nobreak $\dfrac{x^5}{x^2}$\quad（3）\nobreak $x^2\times x^{-1}$\quad（4）\nobreak $\dfrac{a^4}{a}$\quad（5）\nobreak $y^0$\quad（6）\nobreak $\dfrac{10h+5h^2}{h}$\par
\kangaekata かけるときは\R{指数}{し|すう}をたし、わるときはひきます（\R{文字}{も|じ}は0でない\R{数}{かず}）。$a$ は $a^1$ と\R{考}{かんが}えます。\par
\shiki（1）$x^{3+4}=x^7$\quad（2）$x^{5-2}=x^3$\quad（3）$x^{2+(-1)}=x^1=x$\par
\hspace*{2.6\zw}（4）$\dfrac{a^4}{a^1}=a^{4-1}=a^3$\quad（5）0\R{乗}{じょう}は1。$y^0=1$\par
\hspace*{2.6\zw}（6）1つずつわる。$\dfrac{10h+5h^2}{h}=\dfrac{10h}{h}+\dfrac{5h^2}{h}=10+5h$\par
\kotae（1）\nobreak $x^7$\quad（2）\nobreak $x^3$\quad（3）\nobreak $x$\quad（4）\nobreak $a^3$\quad（5）\nobreak 1\quad（6）\nobreak $10+5h$
\end{reidai}

\begin{jibun}
\R{計算}{けい|さん}しましょう。\quad（1）\nobreak $a^2\times a^5$\quad（2）\nobreak $\dfrac{y^6}{y^4}$\quad（3）\nobreak $t^3\times t^{-2}$\quad（4）\nobreak $\dfrac{b^5}{b}$\quad（5）\nobreak $t^0$\quad（6）\nobreak $\dfrac{6t^2+4t}{2t}$
\bkjans{（1）\nobreak $a^7$\quad（2）\nobreak $y^2$\quad（3）\nobreak $t$\quad（4）\nobreak $b^4$\quad（5）\nobreak 1\quad（6）\nobreak $3t+2$}
\end{jibun}


%% ---------------- まちがえやすいところ ----------------
\bksection{tsumazuki}{まちがえやすいところをチェック}

\begin{tsumazuki}{\R{文字}{も|じ}の\R{項}{こう}と\R{数}{かず}の\R{項}{こう}をまとめてしまう}
\bkng{$2x+3=5x$}
\bkok{$2x+3$ は、これ\R{以上}{い|じょう}まとめられない}
$2x$ と3は\R{同類項}{どう|るい|こう}ではありません。\qty{2}{m}と3\R{秒}{びょう}をたせないのと\R{同}{おな}じ。
\end{tsumazuki}

\begin{tsumazuki}{かっこをはずすとき、うしろの\R{項}{こう}にかけわすれる}
\bkng{$3(x+2)=3x+2$}
\bkok{$3(x+2)=3x+6$}
かっこの\R{外}{そと}の3は、\R{中}{なか}の $x$ にも2にもかけます。\R{長方形}{ちょう|ほう|けい}の\R{図}{ず}（\pageref{bk:2-12:kakko}ページ）では、\R{右}{みぎ}の\R{部分}{ぶ|ぶん}の\R{面積}{めん|せき}も $3\times2=6$ です。
\end{tsumazuki}

\begin{tsumazuki}{マイナスのかっこで、うしろの\R{符号}{ふ|ごう}を\R{変}{か}えわすれる}
\bkng{$-(x-2)=-x-2$}
\bkok{$-(x-2)=-x+2$}
$-1$ を、かっこの\R{中}{なか}のどの\R{項}{こう}にもかけます。$(-1)\times(-2)=2$ です。
\end{tsumazuki}

\begin{tsumazuki}{\R{指数}{し|すう}をかけてしまう}
\bkng{$x^2\times x^3=x^6$}
\bkok{$x^2\times x^3=x^5$}
$x$ を\R{何回}{なん|かい}かけたかを\R{数}{かず}えると、$2+3=5$ \R{回}{かい}です。\R{指数}{し|すう}はたします。
\end{tsumazuki}


%% ---------------- 物理・工学 ----------------
\begin{butsuri}
\leavevmode\label{bk:2-12:butsuri}\textbf{\R{単位}{たん|い}も、\R{文字}{も|じ}と\R{同}{おな}じように\R{計算}{けい|さん}する}\quad m や s も、\R{文字}{も|じ}のように\R{計算}{けい|さん}できます（2-10）。
\begin{bkcenter}
\begin{bktable}{colspec={Q[l,m,wd=36mm] Q[l,m,wd=36mm] Q[l,m,wd=42mm]}, rowsep=2.5pt, cells={font=\small}, row{1}={bg=bkPurple!12}}
\R{文字式}{も|じ|しき} & \R{単位}{たん|い} & \R{意味}{い|み} \\
$x\times x=x^2$ & $\mathrm{m}\times\mathrm{m}=\mathrm{m^2}$ & \R{長}{なが}さ×\R{長}{なが}さで\R{面積}{めん|せき} \\
$\dfrac{x}{t}\times t=x$ & $\dfrac{\mathrm{m}}{\mathrm{s}}\times\mathrm{s}=\mathrm{m}$ & \R{速}{はや}さ×\R{時間}{じ|かん}で\R{道}{みち}のり \\
$\dfrac{1}{t}=t^{-1}$ & $\dfrac{1}{\mathrm{s}}=\mathrm{s^{-1}}$ & 1\R{秒}{びょう}あたりの\R{回数}{かい|すう} \\
$3x+2x=5x$ & $\qty{3}{m}+\qty{2}{m}=\qty{5}{m}$ & \R{同}{おな}じ\R{単位}{たん|い}どうしはたせる \\
\end{bktable}
\end{bkcenter}

\noindent\begin{minipage}[c]{0.69\linewidth}
\par\addvspace{3pt}\textbf{コンセントの\R{電気}{でん|き}は\qty{50}{Hz}と\qty{60}{Hz}}\par\noindent コンセントの\R{電気}{でん|き}は、\R{向}{む}きが\R{行}{い}ったり\R{来}{き}たりする\R{波}{なみ}を、\R{東日本}{ひがし|に|ほん}では1\R{秒}{びょう}に50\R{回}{かい}、\R{西日本}{にし|に|ほん}では1\R{秒}{びょう}に60\R{回}{かい}くり\R{返}{かえ}します。この「1\R{秒}{びょう}あたりの\R{回数}{かい|すう}」の\R{単位}{たん|い}がHz（ヘルツ）で、$\mathrm{Hz}=\mathrm{s^{-1}}$ です。Hzも、2-10のNと\R{同}{おな}じように、\R{組}{く}み\R{立}{た}てた\R{単位}{たん|い}に\R{名前}{な|まえ}をつけたものです。\R{大学}{だい|がく}の\R{本}{ほん}では、\unit{m/s} を $\mathrm{m\,s^{-1}}$、\unit{m/s^2} を $\mathrm{m\,s^{-2}}$ と\R{書}{か}きます。$x^{-1}=\dfrac{1}{x}$ のやくそくです。
\end{minipage}\hfill
\begin{minipage}[c]{0.28\linewidth}\centering
\begin{tikzpicture}[x=1mm, y=1mm]
  \filldraw[fill=white, draw=bkGray, line width=0.9pt, rounded corners=2.5mm] (0,0) rectangle (26,30);
  \foreach \yy in {8,22} {\fill[bkInk!80, rounded corners=0.4mm] (8.6,\yy-2.6) rectangle (10.4,\yy+2.6); \fill[bkInk!80, rounded corners=0.4mm] (15.6,\yy-2.6) rectangle (17.4,\yy+2.6);}
  \draw[bkPurple, line width=1.1pt, domain=0:26, samples=60] plot ({\x}, {-9+3.5*sin(\x*360/13)});
  \node[font=\sffamily\footnotesize, anchor=north] at (13,-14) {1\R{秒}{びょう}に50\R{回}{かい}};
\end{tikzpicture}
\end{minipage}
\end{butsuri}

\begin{column}{\R{物理}{ぶつ|り}の\R{有名}{ゆう|めい}な\R{式}{しき}も、×をはぶいて\R{書}{か}く}
「\R{力}{ちから}＝\R{質量}{しつ|りょう}×\R{加速度}{か|そく|ど}」は $F=ma$（\R{加速度}{か|そく|ど} $a$ は、1\R{秒}{びょう}ごとに\R{速}{はや}さが\R{何}{なん}m/sふえるか。\R{単位}{たん|い}は2-10の\unit{m/s^2}）、「\R{電圧}{でん|あつ}＝\R{抵抗}{てい|こう}×\R{電流}{でん|りゅう}」は $V=RI$（2-15）と\R{書}{か}きます。アインシュタインの $E=mc^2$ は、「エネルギー＝\R{質量}{しつ|りょう}×（\R{光}{ひかり}の\R{速}{はや}さ）×（\R{光}{ひかり}の\R{速}{はや}さ）」のことです。$c^2$ は $c\times c$、$mc^2$ は $m\times c\times c$。この\R{章}{しょう}のやくそくを\R{知}{し}っていれば、\R{物理}{ぶつ|り}の\R{本}{ほん}の\R{式}{しき}が\R{読}{よ}めるようになります。
\end{column}

%% ---------------- 章末問題 ----------------
\clearpage
\bkendproblems
\bkproblemgroup{\R{基本}{き|ほん}}{2}

\bkproblem やくそくにしたがって\R{書}{か}きましょう。\par
（1）\nobreak $b\times a\times4$\quad（2）\nobreak $x\times y\times x$\quad（3）\nobreak $m\div2$\quad（4）\nobreak $(-1)\times a$\quad（5）\nobreak $(x-3)\div5$
\bkpans{（1）\nobreak $4ab$\quad（2）\nobreak $x^2y$\quad（3）\nobreak $\dfrac{m}{2}$\quad（4）\nobreak $-a$\quad（5）\nobreak $\dfrac{x-3}{5}$}

\bkproblem まとめましょう。まとめられないときは「まとめられない」と\R{答}{こた}えましょう。\par
（1）\nobreak $6x-x$\quad（2）\nobreak $2a+7-5a-3$\quad（3）\nobreak $3y+y^2$\quad（4）\nobreak $x+4+2x$
\bkpans{（1）\nobreak $5x$\quad（2）\nobreak $-3a+4$\quad（3）\nobreak まとめられない\quad（4）\nobreak $3x+4$}

\bkproblem \R{計算}{けい|さん}しましょう。\quad（1）\nobreak $4(3x-2)$\quad（2）\nobreak $-(5-2y)$\quad（3）\nobreak $(12a+8)\div4$
\bkpans{（1）\nobreak $12x-8$\quad（2）\nobreak $-5+2y$\quad（3）\nobreak $3a+2$}

\bkproblem \R{計算}{けい|さん}しましょう。\quad（1）\nobreak $x^4\times x^3$\quad（2）\nobreak $\dfrac{a^6}{a}$\quad（3）\nobreak $x^{-1}\times x^3$\quad（4）\nobreak $y^0$
\bkpans{（1）\nobreak $x^7$\quad（2）\nobreak $a^5$\quad（3）\nobreak $x^2$\quad（4）\nobreak 1}

\bkproblem \R{次}{つぎ}のうち、\R{正}{ただ}しいものをすべて\R{選}{えら}びましょう。\par
\noindent\hspace*{1\zw}ア\enspace $2x+3x=5x$\hspace{4\zw}イ\enspace $2x+3y=5xy$\par
\noindent\hspace*{1\zw}ウ\enspace $x^2\times x^2=x^4$\hspace{3.4\zw}エ\enspace $-(a+1)=-a+1$
\bkpans{アとウ}

\bkproblem \noindent\begin{minipage}[t]{0.56\linewidth}
\R{右}{みぎ}の\R{長方形}{ちょう|ほう|けい}の\R{面積}{めん|せき}を、かっこを\R{使}{つか}った\R{式}{しき}と、はずした\R{式}{しき}の2\R{通}{とお}りで\R{表}{あらわ}しましょう。
\end{minipage}\hfill
\begin{minipage}[t]{0.34\linewidth}\vspace{1mm}\centering
\begin{tikzpicture}[x=1mm, y=1mm]
  \fill[bkBlue!20] (0,0) rectangle (24,13); \fill[bkOrange!25] (24,0) rectangle (40,13);
  \draw[bkInk, line width=0.9pt] (0,0) rectangle (40,13); \draw[bkInk, line width=0.9pt] (24,0) -- (24,13);
  \bkbrace{(0,14.5)}{(24,14.5)}{$x$}{1}\bkbrace{(24,14.5)}{(40,14.5)}{$4$}{1}
  \node[anchor=east, font=\small] at (-1.5,6.5) {$3$};
\end{tikzpicture}
\end{minipage}
\bkpans{$3(x+4)=3x+12$}

\clearpage
\bkproblemgroup{チャレンジ}{2}

\bkproblem \noindent\begin{minipage}[t]{0.57\linewidth}
たての\R{長}{なが}さが $h$、\R{面積}{めん|せき}が $20h+5h^2$ の\R{長方形}{ちょう|ほう|けい}があります（\R{右}{みぎ}の\R{図}{ず}）。\R{横}{よこ}の\R{長}{なが}さを、$h$ の\R{式}{しき}で\R{表}{あらわ}しましょう。（\R{第}{だい}5\R{巻}{かん}では、このような\R{計算}{けい|さん}で、ボールの\R{速}{はや}さを\R{求}{もと}めます。）
\end{minipage}\hfill
\begin{minipage}[t]{0.34\linewidth}\vspace{1mm}\centering
\begin{tikzpicture}[x=1mm, y=1mm]
  \fill[bkBlue!20] (0,0) rectangle (28,12); \fill[bkOrange!25] (28,0) rectangle (40,12);
  \draw[bkInk, line width=0.9pt] (0,0) rectangle (40,12); \draw[bkInk, line width=0.9pt] (28,0) -- (28,12);
  \node[font=\small] at (14,6) {$20h$}; \node[font=\small] at (34,6) {$5h^2$};
  \node[anchor=east, font=\small] at (-1.5,6) {$h$};
  \bkbrace{(0,-1)}{(40,-1)}{？}{-1}
\end{tikzpicture}
\end{minipage}
\bkpans{$\dfrac{20h+5h^2}{h}=\dfrac{20h}{h}+\dfrac{5h^2}{h}=20+5h$}

\bkproblem $2(x+3)-(x-1)$ を\R{計算}{けい|さん}しましょう。
\bkpans{$2x+6-x+1=x+7$}

\bkproblem $x^2\times x^{-2}$ を、\R{指数}{し|すう}のきまりで\R{計算}{けい|さん}しましょう。また、$x^{-2}=\dfrac{1}{x^2}$ を\R{使}{つか}って\R{計算}{けい|さん}し、2つの\R{答}{こた}えが\R{同}{おな}じになることを\R{確}{たし}かめましょう。
\bkpans{\R{指数}{し|すう}のきまりで $x^{2+(-2)}=x^0=1$。\R{分数}{ぶん|すう}で $x^2\times\dfrac{1}{x^2}=1$。\R{同}{おな}じ}

\begin{furikaeri}
\begin{checklist}
\item ×をはぶき、÷を\R{分数}{ぶん|すう}にして、\R{文字式}{も|じ|しき}を\R{短}{みじか}く\R{書}{か}ける。\hfill{\footnotesize\sffamily\color{bkTheme}→\pageref{bk:2-12:yakusoku}ページ}
\item \R{同類項}{どう|るい|こう}をまとめたり、かっこをはずしたりできる。\hfill{\footnotesize\sffamily\color{bkTheme}→\pageref{bk:2-12:dourui}ページ}
\item \R{指数}{し|すう}のきまりを、\R{文字}{も|じ}を\R{使}{つか}って\R{書}{か}いて\R{使}{つか}える。\hfill{\footnotesize\sffamily\color{bkTheme}→\pageref{bk:2-12:shisu}ページ}
\end{checklist}
\end{furikaeri}

\bkfukidashi{mina}{\R{文字式}{も|じ|しき}を\R{短}{みじか}く\R{書}{か}くやくそくがわかったね。\R{次}{つぎ}の2-13では、\R{文字}{も|じ}に\R{数}{かず}を\R{入}{い}れて、\R{答}{こた}えを\R{出}{だ}してみよう。}
